Burman et al. (2026) provide the calibrated-transport geometry and the transthyretin amyloid cardiomyopathy consumer, but their observed target unit and bounded-overlap population theory differ from this covariate-only target experiment; their results do not supply polynomial-boundary held-out localization for the capped continuation used here. Khan and Tamer (2010) establish the general link between arbitrarily large inverse weights and slower-than-parametric behavior. Ma and Wang (2020) analyze one-sample trimmed inverse-probability weighting and local-polynomial trimming-bias repair under regular variation. Heiler and Kazak (2021) derive stable and self-normalized limits for one-sample inverse-weighted and doubly robust estimators. Orihara, Komukai, and Li (2026) use an exact finite-grid polynomial projection for an overlap-weighted estimand path. None of those results gives the present two-sample fixed-H\"older boundary lower modulus or its critical multiscale logarithm. Armstrong and Koles\'ar (2021) solve a conditional-on-realized-design, one-sample unconfounded ATE optimal-recovery problem under smoothness or shape restrictions, allowing limited or absent overlap, and obtain finite-sample Gaussian-known-variance optimal estimators and intervals together with feasible asymptotic procedures. Their experiment has no independent target sample, no triangular array with \(f_S(x)\asymp x^a\) and positive-boundary \(f_T\), and no \(a=1\) multiscale lower modulus. Conversely, this note proves no fixed-design finite-sample optimal estimator or optimal interval, so neither theorem implies the other. Donoho (1994) supplies the modulus-of-continuity blueprint, while Cai and Low (2004) provide the ordered-modulus language for the retained nonadaptation theorem.